\documentclass[aspectratio=169,10pt]{beamer}

% ------------------------------------------------------------------ theme
\usetheme[progressbar=frametitle,block=fill]{metropolis}
\usefonttheme{professionalfonts}
\usepackage[scaled=0.92]{helvet}
\usepackage{courier}
\renewcommand{\familydefault}{\sfdefault}

\usepackage[T1]{fontenc}
\usepackage[utf8]{inputenc}
\usepackage{amsmath,amssymb}
\usepackage{booktabs}
\usepackage{listings}
\usepackage{tcolorbox}
\usepackage{appendixnumberbeamer}
\tcbuselibrary{skins,breakable}

% ------------------------------------------------------------------ colour
\definecolor{accent}{HTML}{EB6834}
\definecolor{cool}{HTML}{2A78D6}
\definecolor{ink}{HTML}{1B1B1B}
\definecolor{paper}{HTML}{FCFCFB}
\definecolor{muted}{HTML}{898781}
\definecolor{good}{HTML}{1C7C54}

\setbeamercolor{normal text}{fg=ink,bg=paper}
\setbeamercolor{alerted text}{fg=accent}
\setbeamercolor{example text}{fg=cool}
\setbeamercolor{frametitle}{bg=ink,fg=paper}
\setbeamercolor{progress bar}{fg=accent,bg=muted!30}
\setbeamercolor{title separator}{fg=accent}
\setbeamercolor{structure}{fg=accent}
\setbeamertemplate{itemize items}[circle]
\setbeamertemplate{section in toc}[sections numbered]
\setbeamerfont{frametitle}{size=\large}
\metroset{numbering=fraction}

% ------------------------------------------------------------------ navigation
% Clickable section strip in the footer, plus a link back to the contents.
\usepackage{hyperref}
\hypersetup{colorlinks=true, linkcolor=accent, urlcolor=cool,
            pdfstartview={Fit}, bookmarksopen=true}

\setbeamercolor{section in head/foot}{fg=muted,bg=paper}
\setbeamercolor{section in head/foot shaded}{fg=muted!60,bg=paper}
\setbeamerfont{section in head/foot}{size=\tiny}

\newcommand{\navbutton}[2]{%
  \hyperlink{#1}{\tiny\color{muted}\textsf{#2}}}

\newcommand{\navitem}[3]{%
  \hyperlink{#1}{%
    \ifnum\insertsectionnumber=#2\color{accent}\else\color{muted!75}\fi
    \textsf{#3}}}

\newcommand{\secnav}{%
  \navitem{sec1}{1}{1\,Components}\hspace{0.7em}%
  \navitem{sec2}{2}{2\,Extract}\hspace{0.7em}%
  \navitem{sec3}{3}{3\,Transform}\hspace{0.7em}%
  \navitem{sec4}{4}{4\,Load}\hspace{0.7em}%
  \navitem{sec5}{5}{5\,One\,function}\hspace{0.7em}%
  \navitem{sec6}{6}{6\,Protect\,it}}

\makeatletter
\setbeamertemplate{footline}{%
  \hbox{%
    \begin{beamercolorbox}[wd=\paperwidth,ht=2.8ex,dp=1.6ex,%
      leftskip=1.1em,rightskip=1.1em]{}%
      \tiny%
      \navbutton{toc}{Contents}\hspace{1.2em}%
      \secnav%
      \hfill\color{muted}\insertframenumber/\inserttotalframenumber%
    \end{beamercolorbox}%
  }%
}
\makeatother

% Every section opens with its own slide, which doubles as a link target.
\AtBeginSection[]{%
  \begin{frame}[plain,noframenumbering]
    \hypertarget{sec\arabic{section}}{}%
    \sectionpage
  \end{frame}%
}
\metroset{numbering=none}

% A prompt that asks the room to answer before the next click.
\newtcolorbox{ask}[1][]{
  enhanced, breakable, colback=good!6, colframe=good,
  boxrule=0pt, leftrule=2.5pt, arc=1pt, left=6pt, right=6pt, top=4pt, bottom=4pt,
  fonttitle=\bfseries\footnotesize, title={Before the next click}, #1}


% ------------------------------------------------------------------ code
\lstdefinestyle{py}{
  language=Python,
  basicstyle=\ttfamily\footnotesize,
  keywordstyle=\color{accent}\bfseries,
  stringstyle=\color{good},
  commentstyle=\color{muted}\itshape,
  showstringspaces=false,
  breaklines=true,
  frame=none,
  aboveskip=4pt,
  belowskip=2pt,
  numbers=none,
  xleftmargin=2pt,
}
\lstset{style=py}

% Inline code. (\lstinline cannot be used inside a macro argument,
% so inline snippets use a plain typewriter face instead.)
\newcommand{\py}[1]{{\ttfamily\footnotesize\color{ink}#1}}
\newcommand{\Out}{\vspace{-2pt}\textcolor{muted}{\scriptsize\textsf{Output}}\vspace{-4pt}}

% ------------------------------------------------------------------ boxes
\newtcolorbox{defn}[1][]{
  enhanced, breakable, colback=cool!5, colframe=cool,
  boxrule=0pt, leftrule=2.5pt, arc=1pt, left=6pt, right=6pt, top=4pt, bottom=4pt,
  fonttitle=\bfseries\footnotesize, title={Definition}, #1}

\newtcolorbox{exmp}[1][]{
  enhanced, breakable, colback=accent!5, colframe=accent,
  boxrule=0pt, leftrule=2.5pt, arc=1pt, left=6pt, right=6pt, top=4pt, bottom=4pt,
  fonttitle=\bfseries\footnotesize, title={Example}, #1}

\newtcolorbox{aside}[1][]{
  enhanced, breakable, colback=muted!8, colframe=muted,
  boxrule=0pt, leftrule=2.5pt, arc=1pt, left=6pt, right=6pt, top=4pt, bottom=4pt, #1}

% ------------------------------------------------------------------ title
\title{Building a Simple ETL Pipeline}
\subtitle{Components of data pipelines \textperiodcentered{} Extract, Transform, Load
          \textperiodcentered{} Making it run again tomorrow}
\date{}
\author{Moussa Doumbia, Ph.D.\\[2pt]
        {\small Expert in DS, ML \& AI}\\[1pt]
        {\small Department of Mathematics}\\[1pt]
        {\small Howard University}}
\institute{MATH 153 \textemdash{} Data Wrangling / Data Engineering}

\newtcolorbox{prac}[1][]{
  enhanced, colback=good!6, colframe=good,
  boxrule=0pt, leftrule=2.5pt, arc=1pt, left=6pt, right=6pt, top=4pt, bottom=4pt,
  coltext=ink, fonttitle=\bfseries\footnotesize, title={In practice}, #1}

\newtcolorbox{watch}[1][]{
  enhanced, colback=accent!7, colframe=accent,
  boxrule=0pt, leftrule=2.5pt, arc=1pt, left=6pt, right=6pt, top=4pt, bottom=4pt,
  coltext=ink, fonttitle=\bfseries\footnotesize, title={Watch out}, #1}
\begin{document}

\maketitle

% ==================================================================
\begin{frame}{Today}
  \hypertarget{toc}{}%
  \tableofcontents
\end{frame}

% ==================================================================
\section{What a pipeline is}

\begin{frame}{The job}
  \begin{itemize}[<+->]
    \item The Bison Caf\'e sends us \textbf{one sales file per day}.
    \item We have three days.
    \item We want a clean table, and a daily revenue report.
  \end{itemize}

  \pause\vfill
  \begin{aside}
    \small
    And we want it to work again tomorrow, \emph{without us}.
  \end{aside}
\end{frame}

% ==================================================================
\begin{frame}{Definition}
  \begin{defn}
    A \textbf{data pipeline} is a series of steps that moves data from where it \emph{is} to where it
    is \emph{useful}, doing the same thing every time it runs.
  \end{defn}

  \pause
  \begin{itemize}[<+->]
    \item The important words are \alert{every time it runs}.
    \item A notebook you run once by hand is an \textbf{analysis}.
    \item A pipeline runs tomorrow, and next Tuesday, without you.
  \end{itemize}
\end{frame}

% ==================================================================
\begin{frame}{The five components}
  \small
  \begin{tabular}{@{}p{0.4cm}p{2.4cm}p{4.3cm}p{5.1cm}@{}}
    \toprule
    \textbf{\#} & \textbf{Component} & \textbf{What it means} & \textbf{In our example} \\
    \midrule
    \onslide<+->{1} & \onslide<.->{\textbf{Source}}
       & \onslide<.->{where the data comes from} & \onslide<.->{3 daily CSVs from the till} \\[2pt]
    \onslide<+->{2} & \onslide<.->{\textbf{Extract}}
       & \onslide<.->{read it in} & \onslide<.->{\py{pd.read\_csv}} \\[2pt]
    \onslide<+->{3} & \onslide<.->{\textbf{Transform}}
       & \onslide<.->{make it usable} & \onslide<.->{fix prices, fix names, join the menu} \\[2pt]
    \onslide<+->{4} & \onslide<.->{\textbf{Load}}
       & \onslide<.->{put it where people will use it} & \onslide<.->{a SQLite table + a summary CSV} \\[2pt]
    \onslide<+->{5} & \onslide<.->{\textbf{Schedule}}
       & \onslide<.->{decide when it runs} & \onslide<.->{every morning at 6 a.m.} \\
    \bottomrule
  \end{tabular}

  \onslide<6->{\vfill
  \begin{aside}
    \small
    Steps 2, 3 and 4 are so common they have a name of their own: \alert{ETL}.
  \end{aside}}
\end{frame}

% ==================================================================
\begin{frame}[fragile]{E \textemdash{} T \textemdash{} L}
\begin{lstlisting}[language=,basicstyle=\ttfamily\scriptsize]
 SOURCE      EXTRACT     TRANSFORM      LOAD       SCHEDULE
 3 CSVs  --> read    --> clean+join --> cafe.db <-- daily
 + menu                                 summary     6 a.m.
\end{lstlisting}

  \pause
  \begin{itemize}[<+->]
    \item \textbf{Extract} \textemdash{} get the raw data, \emph{without changing it}.
    \item \textbf{Transform} \textemdash{} clean it, reshape it, combine it, calculate from it.
    \item \textbf{Load} \textemdash{} write the result somewhere other people can read.
  \end{itemize}
\end{frame}

% ==================================================================
\begin{frame}{ETL or ELT?}
  \begin{itemize}[<+->]
    \item \textbf{ETL} \textemdash{} transform first, then load the clean result.
    \item \textbf{ELT} \textemdash{} load the raw data first, transform it \emph{inside} the database.
  \end{itemize}

  \pause\vfill
  \begin{aside}
    \small
    Same three letters, different order. ELT suits very large data and powerful warehouses.
    ETL suits what we are doing today.
  \end{aside}
\end{frame}

% ==================================================================
\section{Extract}

\begin{frame}[fragile]{Look at the data first}
\begin{lstlisting}[language=]
sale_id,sale_date,sale_time,item,qty,unit_price
S-5001,2027-03-02,14:32,FRUIT CUP,2,$4.00
S-5002,2027-03-02,15:22,Turkey Wrap,2,$7.50
S-5004,2027-03-02,11:58,latte,3,$4.25
\end{lstlisting}

  \pause
  \begin{ask}
    \small
    Two things in this file are already wrong. What are they?
  \end{ask}
\end{frame}

% ==================================================================
\begin{frame}{Two problems, both normal}
  \begin{itemize}[<+->]
    \item \py{unit\_price} \textbf{is text, not a number}. It says \py{\$4.25}.
    \item You cannot multiply by \py{\$4.25}.
    \item \py{item} \textbf{is spelled inconsistently} \textemdash{} \py{FRUIT CUP}, \py{latte},
          \py{"  Bagel "}.
    \item The till lets staff type it by hand.
  \end{itemize}

  \pause\vfill
  \begin{aside}
    \small
    Fixing those two things is what ``Transform'' will mean today.
  \end{aside}
\end{frame}

% ==================================================================
\begin{frame}[fragile]{Extract: read, and change nothing}
\begin{lstlisting}
def extract(data_dir):
    """Read every daily sales file and stack them into one DataFrame."""
    files = sorted(data_dir.glob("sales_2027-*.csv"))
    frames = [pd.read_csv(f) for f in files]
    return pd.concat(frames, ignore_index=True)
\end{lstlisting}
  \Out
\begin{lstlisting}[language=]
107 rows from 3 files
\end{lstlisting}

  \pause
  \begin{aside}
    \small
    \textbf{Extract is deliberately boring.} Keeping the raw data untouched means that when something
    goes wrong later, you know where to start again.
  \end{aside}
\end{frame}

% ==================================================================
\begin{frame}[fragile]{What \texttt{glob} is for}
\begin{lstlisting}
files = ["sales_2027-03-02.csv", "sales_2027-03-03.csv", ...]  # breaks Friday
files = sorted(data_dir.glob("sales_2027-*.csv"))              # just works
\end{lstlisting}

  \pause
  \begin{itemize}[<+->]
    \item \py{glob} finds files by \alert{pattern} instead of by name.
    \item So tomorrow's file is picked up without editing the code.
    \item The pattern also \emph{excludes} \py{menu.csv}, which lives in the same folder.
    \item \py{sorted} matters: \py{glob} returns files in filesystem order, not date order.
  \end{itemize}
\end{frame}

% ==================================================================
\section{Transform}

\begin{frame}[fragile]{Step 1 \textemdash{} make the price a number}
\begin{lstlisting}
df["unit_price"] = (df["unit_price"]
                      .str.replace("$", "", regex=False)
                      .astype(float))
\end{lstlisting}
  \Out
\begin{lstlisting}[language=]
before: $4.00   (a str)
after : 4.0     (a float64)
\end{lstlisting}

  \pause
  \begin{aside}
    \small
    \py{regex=False} because \py{\$} means ``end of string'' in a regular expression. Here we want the
    literal dollar sign.
  \end{aside}
\end{frame}

% ==================================================================
\begin{frame}[fragile]{Step 2 \textemdash{} make the date a real date}
\begin{lstlisting}
df["sale_date"] = pd.to_datetime(df["sale_date"], format="%Y-%m-%d")
\end{lstlisting}

  \pause
  \begin{itemize}[<+->]
    \item A date stored as text is just a string.
    \item You cannot sort it reliably, subtract it, or ask what day of the week it was.
    \item Convert it \emph{once}, here, and every later step gets a real date.
  \end{itemize}
\end{frame}

% ==================================================================
\begin{frame}[fragile]{Step 3 \textemdash{} clean the item names}
\begin{lstlisting}
df["item"] = df["item"].str.strip().str.title()
\end{lstlisting}
  \Out
\begin{lstlisting}[language=]
distinct spellings before: 29
distinct item names after:  8
\end{lstlisting}

  \pause
  \begin{itemize}[<+->]
    \item \py{.str.strip()} removes spaces at the ends \textemdash{} \py{"  Latte "} becomes
          \py{"Latte"}.
    \item \py{.str.title()} capitalises each word \textemdash{} \py{"LATTE"} and \py{"latte"} both
          become \py{"Latte"}.
  \end{itemize}
\end{frame}

% ==================================================================
\begin{frame}[fragile]{Step 4 \textemdash{} join the menu, then calculate}
\begin{lstlisting}
df = df.merge(menu, on="item", how="left")
df["revenue"] = (df["qty"] * df["unit_price"]).round(2)
\end{lstlisting}

  \pause
  \begin{ask}
    \small
    Why did step 3 have to come \emph{before} step 4?
  \end{ask}

  \pause
  \begin{watch}
    \small
    \py{"latte"} would not have matched \py{"Latte"} in the menu. That sale would have had no
    category, and its revenue would have \alert{quietly vanished} from the report.
  \end{watch}
\end{frame}

% ==================================================================
\begin{frame}{Two checks after every join}
  \begin{itemize}[<+->]
    \item \textbf{Did the row count stay the same?} 107 in, 107 out.
    \item If it grew, the reference table has duplicate rows and every sale got counted twice.
    \item \textbf{Did anything fail to match?} Zero unmatched here.
    \item A blank category means a spelling your cleaning step missed.
  \end{itemize}

  \pause\vfill
  \begin{prac}
    \small
    Print both numbers. A join that silently doubled your data looks exactly like one that worked.
  \end{prac}
\end{frame}

% ==================================================================
\section{Load}

\begin{frame}{Two outputs, two audiences}
  \begin{itemize}[<+->]
    \item A \textbf{database table} \textemdash{} for anyone who wants to query it.
    \item A \textbf{summary CSV} \textemdash{} for the manager who just wants the daily numbers.
  \end{itemize}

  \pause\vfill
  \begin{aside}
    \small
    ``Load'' does not mean one place. It means wherever the people who need the answer will actually
    look.
  \end{aside}
\end{frame}

% ==================================================================
\begin{frame}[fragile]{Load}
\begin{lstlisting}
con = sqlite3.connect("cafe.db")
out.to_sql("sales", con, if_exists="replace", index=False)
con.close()

summary.to_csv("daily_summary.csv", index=False)
\end{lstlisting}

  \pause
  \begin{aside}
    \small
    Dates go into SQLite as \py{'YYYY-MM-DD'} text, because SQLite has no date type \textemdash{} and
    in that format, alphabetical order \emph{is} chronological order.
  \end{aside}
\end{frame}

% ==================================================================
\begin{frame}{The result}
  \small
  \begin{tabular}{@{}llrr@{}}
    \toprule
    \textbf{sale\_date} & \textbf{category} & \textbf{items\_sold} & \textbf{revenue} \\
    \midrule
    \onslide<+->{2027-03-02} & \onslide<.->{Drink} & \onslide<.->{35} & \onslide<.->{\$119.75} \\
    \onslide<+->{2027-03-02} & \onslide<.->{Food}  & \onslide<.->{40} & \onslide<.->{\$222.85} \\
    \onslide<+->{2027-03-03} & \onslide<.->{Drink} & \onslide<.->{48} & \onslide<.->{\$162.80} \\
    \onslide<+->{2027-03-03} & \onslide<.->{Food}  & \onslide<.->{29} & \onslide<.->{\$112.40} \\
    \onslide<+->{2027-03-04} & \onslide<.->{Drink} & \onslide<.->{39} & \onslide<.->{\$124.30} \\
    \onslide<+->{2027-03-04} & \onslide<.->{Food}  & \onslide<.->{29} & \onslide<.->{\$126.25} \\
    \midrule
    \onslide<+->{\textbf{Total}} & & & \onslide<.->{\textbf{\$868.35}} \\
    \bottomrule
  \end{tabular}

  \onslide<8->{\vfill
  \begin{aside}
    \small
    107 sales over three days, from two messy columns and a join.
  \end{aside}}
\end{frame}

% ==================================================================
\section{The whole thing, in one function}

\begin{frame}[fragile]{\texttt{run\_pipeline()}}
\begin{lstlisting}[basicstyle=\ttfamily\scriptsize]
def run_pipeline(data_dir=DATA, db_path="cafe.db", csv_path="daily_summary.csv"):
    """Extract -> Transform -> Load. Run this every morning."""
    menu  = pd.read_csv(data_dir / "menu.csv")

    raw   = extract(data_dir)                       # E
    clean = transform(raw, menu)                    # T

    assert clean["category"].notna().all(), "some items did not match the menu"
    assert (clean["revenue"] > 0).all(),    "a sale has zero or negative revenue"

    summary = (clean
               .assign(day=clean["sale_date"].dt.date)
               .groupby(["day", "category"], as_index=False)
               .agg(items_sold=("qty", "sum"), revenue=("revenue", "sum"))
               .rename(columns={"day": "sale_date"}))

    load(clean, summary, db_path, csv_path)         # L
    return clean, summary
\end{lstlisting}

  \pause
  \begin{aside}
    \small
    Twelve lines. \alert{This is the deliverable} \textemdash{} not the notebook, the function.
  \end{aside}
\end{frame}

% ==================================================================
\begin{frame}{Why a function, and not just cells?}
  \begin{itemize}[<+->]
    \item \textbf{It has a name.}
    \item You can say ``run \py{run\_pipeline()}'' instead of ``run cells 4 to 19 but not cell 12''.
    \item \textbf{It takes arguments.}
    \item Point it at a different folder and it works on different data.
    \item \textbf{It can be scheduled.}
    \item A computer can call a function at 6 a.m. It cannot click through a notebook.
  \end{itemize}
\end{frame}

% ==================================================================
\section{Protecting it}

\begin{frame}[fragile]{Run it twice}
\begin{lstlisting}
run_pipeline()
run_pipeline()          # on purpose
\end{lstlisting}

  \pause
  \begin{ask}
    \small
    How many rows are in the table now?
  \end{ask}

  \pause
\begin{lstlisting}[language=]
rows in the sales table: 107        <-- not 214
\end{lstlisting}

  \pause
  \begin{aside}
    \small
    Running it twice gave the same answer as running it once. That property has a name: the pipeline
    is \alert{idempotent}.
  \end{aside}
\end{frame}

% ==================================================================
\begin{frame}[fragile]{One keyword makes the difference}
\begin{lstlisting}
out.to_sql("sales", con, if_exists="replace", index=False)   # yes
out.to_sql("sales", con, if_exists="append",  index=False)   # NO
\end{lstlisting}

  \pause
  \begin{watch}
    \small
    With \py{append}, every run adds another copy. Run it three times and the caf\'e has sold
    everything three times.

    \medskip
    This is the \alert{single most common bug in a first pipeline}.
  \end{watch}
\end{frame}

% ==================================================================
\begin{frame}[fragile]{A log line per step}
\begin{lstlisting}[language=]
  EXTRACT    107 rows read
  TRANSFORM  107 rows cleaned, 8 distinct items
  LOAD       107 rows written, revenue $868.35
\end{lstlisting}

  \pause
  \begin{itemize}[<+->]
    \item Three \py{print} statements.
    \item Now when the job runs at 6 a.m.\ and something looks odd \ldots{}
    \item \ldots{} you can see \alert{which step} the rows disappeared at, instead of guessing.
  \end{itemize}
\end{frame}

% ==================================================================
\begin{frame}[fragile]{A check before you publish}
\begin{lstlisting}
assert clean["category"].notna().all(), "some items did not match the menu"
assert (clean["revenue"] > 0).all(),    "a sale has zero or negative revenue"
\end{lstlisting}

  \pause
  \begin{itemize}[<+->]
    \item An \textbf{assertion} is a sentence that must be true.
    \item If it is not, Python stops immediately and tells you why.
  \end{itemize}

  \pause
  \begin{ask}
    \small
    Where in the pipeline should these two lines go?
  \end{ask}
\end{frame}

% ==================================================================
\begin{frame}{Between Transform and Load}
  \begin{itemize}[<+->]
    \item That is the \alert{last moment} before the numbers reach somebody else.
    \item If a check fails, the pipeline stops \emph{before} writing to the database.
    \item The manager's report is still yesterday's.
    \item Which is correct \textemdash{} and much better than today's report being wrong.
  \end{itemize}

  \pause\vfill
  \begin{prac}
    \small
    A pipeline that publishes a wrong number is worse than a pipeline that publishes nothing.
  \end{prac}
\end{frame}

% ==================================================================
\begin{frame}[fragile]{Watch a check fail}
\begin{lstlisting}
broken = clean.copy()
broken.loc[0, "revenue"] = -5.0          # pretend the transform has a bug

assert (broken["revenue"] > 0).all(), "a sale has zero or negative revenue"
\end{lstlisting}
  \Out
\begin{lstlisting}[language=]
Pipeline stopped: a sale has zero or negative revenue
\end{lstlisting}

  \pause
  \begin{aside}
    \small
    Nothing was written to the database. The bad number never left the building.
  \end{aside}
\end{frame}

% ==================================================================
\begin{frame}[fragile]{The same summary, in SQL}
\begin{lstlisting}[language=SQL,basicstyle=\ttfamily\footnotesize,
                   keywordstyle=\color{accent}\bfseries]
SELECT sale_date, category,
       SUM(qty)               AS items_sold,
       ROUND(SUM(revenue), 2) AS revenue
FROM   sales
GROUP  BY sale_date, category
ORDER  BY sale_date, category;
\end{lstlisting}

  \pause
  \begin{aside}
    \small
    Once the data is in the database, the report is one query. Same answer, different language.
  \end{aside}
\end{frame}

% ==================================================================
\begin{frame}{pandas or SQL?}
  \small
  \begin{tabular}{@{}p{6.2cm}p{6.2cm}@{}}
    \toprule
    \textbf{Use pandas when} & \textbf{Use SQL when} \\
    \midrule
    \onslide<+->{reading files and cleaning text}
      & \onslide<.->{the data is already in a database} \\[3pt]
    \onslide<+->{the step is awkward to express in SQL}
      & \onslide<.->{the data is too big for memory} \\[3pt]
    \onslide<+->{you are still exploring}
      & \onslide<.->{the question is a standard group-and-total} \\
    \bottomrule
  \end{tabular}

  \onslide<4->{\vfill
  \begin{aside}
    \small
    Most real pipelines use both: pandas to clean and load, SQL to query afterwards. That is exactly
    what we did.
  \end{aside}}
\end{frame}

% ==================================================================
\section*{Summary}

\begin{frame}{Five things to remember}
  \begin{enumerate}[<+->]
    \item \textbf{Extract changes nothing.} Keep the raw data raw, so you can always start over.
    \item \textbf{Clean before you join.} \py{"latte"} does not match \py{"Latte"}.
    \item \textbf{Check the row count after every join.} Same number in, same number out.
    \item \textbf{Make it runnable twice.} \py{if\_exists="replace"}, not \py{"append"}.
    \item \textbf{Put the checks between Transform and Load.}
  \end{enumerate}
\end{frame}

% ==================================================================
\begin{frame}{The one-sentence version}
  \vfill
  \begin{center}
    {\large A pipeline is just \alert{three functions and a schedule}.}
  \end{center}
  \pause
  \vfill
  \begin{center}
    {\large The difference between a pipeline and a notebook\\[4pt]
     is that \alert{somebody else can run it}.}
  \end{center}
  \vfill
\end{frame}

% ==================================================================
\begin{frame}{In-class activity}
  \begin{exmp}[title=Founders Library checkouts]
    \small
    Three daily files. \py{loan\_days} arrives as \py{"14 days"}. \py{patron\_type} is typed
    inconsistently. Build the pipeline.
  \end{exmp}

  \pause
  \begin{itemize}[<+->]
    \item Six tasks, about 35 minutes, in pairs.
    \item Same five components \textemdash{} only the data is different.
    \item Notebook: \py{ETL\_Activity.ipynb}
  \end{itemize}

  \pause\vfill
  \begin{prac}
    \small
    Expected: \textbf{82 rows}, \textbf{16 spellings become 4}, and the table still says 82 after you
    run it twice.
  \end{prac}
\end{frame}

\end{document}
